\bookchapter{Kapitel 8: Kommunikation: SMS, Telefon und soziale Medien}

\chapterintro{Kommunikation ist die ursprünglichste und wichtigste Aufgabe eines Smartphones. Aus einfachen Anrufen und SMS ist inzwischen ein komplexes Ökosystem geworden, das Instant Messaging, soziale Netzwerke und E-Mail umfasst.}

\booksection{Erweiterte Funktionen der Telefon-App}

Die Telefon-App moderner Smartphones ist längst mehr als ein einfaches Wählpad: Sie integriert leistungsstarke Werkzeuge wie Anrufaufzeichnung, Spam-Blockierung und intelligente Erkennung.

\booksubsection{Gesprächsaufzeichnung (rechtliche Vorgaben beachten):}

\begin{itemize}
  \item Zweck: Wichtige Gesprächsinhalte festhalten – für Protokolle, Interviews, Kundenzusagen oder zur Beweissicherung bei Streitigkeiten.
  \item Rechtlicher Hinweis: In China ist die Verwendung einer Gesprächsaufzeichnung als Beweis an die Einwilligung der anderen Seite gebunden oder nur zulässig, wenn sie der Wahrung der eigenen berechtigten Interessen dient.
  \item Einseitige Aufzeichnung vs. beidseitige Einwilligung: Rein rechtlich sind Aufnahmen ohne Einwilligung der anderen Seite nicht automatisch ungültig. Nach einschlägigen Auslegungen des Obersten Volksgerichtshofs können sie als Beweis vorgelegt werden, sofern die Art der Aufnahme nicht schwerwiegend in Rechte Dritter eingreift, keine gesetzlichen Verbote verletzt und nicht gegen die guten Sitten verstößt. In einigen Provinzen sind die Gerichte im Zivilprozess bei „heimlichen“ Aufnahmen jedoch strenger; in Strafsachen wird in der Regel eine strengere Prüfung der Rechtmäßigkeit verlangt.
  \item Regionale Unterschiede: Gerichte in Peking, Shanghai, Guangdong und anderen Regionen gehen mit einseitigen Aufnahmen als Zivilbeweis eher großzügig um. Aufnahmen, die Geschäftsgeheimnisse oder Privatsphäre betreffen, erfordern dagegen überall besondere Vorsicht und können wegen Verletzung des Persönlichkeitsrechts ausgeschlossen werden.
  \item So geht's: Unter Android gibt es im Anrufbildschirm meist eine deutliche „Aufnehmen“-Taste. iOS bietet keine native Gesprächsaufzeichnung; hier brauchen Sie eine Drittanbieter-App oder ein Bluetooth-Headset mit Aufnahmefunktion.
\end{itemize}

\begin{tcbwarn}{Rote Linien bei der Nutzung von Aufnahmen}
\begin{itemize}
  \item Aufnahmen dürfen nicht zur Erpressung, Verbreitung oder öffentlichen Beleidigung anderer missbraucht werden.
  \item Inhalte, die die Privatsphäre Dritter betreffen (z. B. versehentlich mitgezeichnete Gespräche), dürfen nicht einfach weitergegeben werden.
  \item In wichtigen Fällen empfiehlt es sich, zu Gesprächsbeginn mündlich mitzuteilen: „Dieses Gespräch wird aufgezeichnet.“ Das ist nicht nur rechtssicher, sondern stärkt auch die Beweiskraft.
\end{itemize}

\end{tcbwarn}

\booksubsection{Belästigungsanrufe im Detail:}

\begin{itemize}
  \item Funktionsweise: Die Telefon-App gleicht eingehende Anrufe online mit riesigen Nummern-Datenbanken ab und zeigt „Spamanruf“, „Betrugsanruf“ oder den Firmennamen an.
  \item So aktivieren Sie es: Android: Telefon-App > Einstellungen > Spam-Schutz. iOS: Einstellungen > Telefon > Anrufe blockieren \& identifizieren. Detaillierte Schaltwege und empfohlene Optionen je Marke finden Sie in der Markenübersicht in Kapitel 13 „Praktische Systemwerkzeuge“.
  \item Die wichtigsten Markierungsplattformen: In China stammen die Rufnummern-Markierungen hauptsächlich von drei Datendiensten – Tencent Mobile Manager, 360 Mobile Guard und Dianhuabang (Tapir). Auch die Anzeige des „Inhabernamens“ und die Funktion „Belästigung blockieren“ der Mobilfunkanbieter (China Mobile/China Unicom/China Telecom) greifen oft auf diese Daten zu. Die Markierungsdatenbanken der Handyhersteller (Huawei, Xiaomi, OPPO, vivo) arbeiten meist mit den genannten Diensten zusammen.
  \item Was tun bei einer falschen Markierung? Wenn deine normale Arbeitsnummer fälschlich als „Essenslieferung“, „Paketdienst“ oder „Belästigungsanruf“ markiert wurde, kannst du über folgende Kanäle Widerspruch einlegen: Dianhuabang: Ruf die Beschwerdeplattform für Rufnummern auf (haoma.tk.cn) und reiche einen Identitätsnachweis ein.
  \item 360: Über die Funktion „Nummern-Beschwerde“ in der 360 Mobile Guard App einreichen.
  \item Tencent: Über die Seite „Nummern-Beschwerde“ im Mini-Programm „Tencent Mobile Manager“.
  \item Die drei großen Anbieter: China-Mobile-Kunden senden „KTFSR“ per SMS an 10086, um den Anti-Spam-Dienst zu aktivieren. Telecom-Kunden folgen dem Konto „Tianyi Anti-Spam“; Unicom-Kunden richten es im „Unicom Mobile Manager“ ein.
\end{itemize}

Fortgeschrittener Schutz: Mit der „Hochfrequenz-Anrufsperre“ des Anbieters (z. B. Chinas Mobile „Hochfrequenz-Anrufbegrenzung“) lassen sich automatisch KI-Sprachanrufe blockieren.

Typische Blockierfunktionen des Smartphones: Anruferkennung / intelligente Blockierung: Spam, Betrug und Werbeanrufe werden anhand der Nummern-Datenbank automatisch markiert und blockiert.

Schwarze Liste / Weiße Liste: Bekannte Spam-Nummern manuell blockieren; Kontakte, Paket- und Lieferdienste auf die weiße Liste setzen, damit sie nicht fälschlich blockiert werden.

SMS-Keyword-Filter: Nachrichten mit Begriffen wie „Gewinn“, „Kredit“ oder „Link anklicken“ landen automatisch im Spam-Ordner.

iOS: Unbekannte Anrufe stumm schalten: Anrufe außerhalb der Kontakte werden nur stumm in die Mailbox geleitet und stören nicht. Gut für alle, die viele Spam-Anrufe bekommen – achten Sie aber darauf, dass wichtige Anrufe von Unbekannten nicht untergehen.

Nicht-stören-Modus: Nachts oder im Meeting nur Anrufe aus den Kontakten zulassen, kombiniert mit Zeitplänen – die Schlafqualität verbessert sich sofort.

Welche Apps ergänzen sinnvoll? Zuerst den integrierten System-Schutz nutzen, dann den nationalen Anti-Betrugs-Center und den Anbieter-Schutz ergänzen. Als Cleaner sind leichte Werkzeuge wie Qingzhuo möglich – aber pro Kategorie (Sicherheit/Reinigung) reicht eine App. Details in Kapitel 13.

\booksubsection{Was sind VoLTE/VoNR und wie aktiviert man sie?}

\begin{itemize}
  \item VoLTE (Voice over LTE): Sprachtelefonie über das 4G-Netz. Beim Telefonieren gibt es keinen Rückfall auf 2G/3G mehr – Sie können gleichzeitig telefonieren und online sein, die Sprachqualität ist klarer und die Verbindung kommt schneller zustande.
  \item VoNR (Voice over NR): Sprachtelefonie über 5G – die nächste Generation nach VoLTE. Niedrigere Latenz, bessere Klangqualität und HD-Videotelefonie. Voraussetzung: 5G-Handy + 5G-Tarif + lokaler 5G-Empfang.
  \item So aktivieren Sie es: iOS: „Einstellungen > Mobilfunk > Mobilfunkdaten-Optionen > Sprach \& Daten“ und „VoLTE“ einschalten. Android: „Einstellungen > Mobilfunk > Mobilfunkdaten > Erweitert“ und den Schalter „VoLTE-HD-Anrufe“ finden. Anbieter schalten VoLTE meist standardmäßig ein. Fällt beim Telefonieren die Internetverbindung weg, ist VoLTE oft nicht aktiviert.
\end{itemize}

\begin{tcbinfo}{Kurz erklärt}
Wird das Internet beim Telefonieren plötzlich langsam oder bricht ab, sind VoLTE/VoNR meist nicht aktiviert – das Gerät fällt auf 2G/3G zurück. Nach dem Einschalten können Sie problemlos „telefonieren und dabei Videos schauen“.

\end{tcbinfo}

\booksection{SMS und iMessage/RCS}

Die klassische SMS ist nicht tot – sie bleibt ein verlässlicher Kanal für Bestätigungscodes, Bankbenachrichtigungen und offizielle Mitteilungen.

\begin{itemize}
  \item iMessage: Nachrichten zwischen Apple-Geräten über das Internet, angezeigt als blaue Sprechblasen. Unterstützt hochwertige Fotos und Videos, Lesebestätigungen usw., komplett kostenlos und Ende-zu-Ende-verschlüsselt.
  \item RCS: Das nächste Messaging-Protokoll, getrieben von Anbietern weltweit und Google. Funktional ähnlich wie iMessage oder WeChat, aber basierend auf der Telefonnummer.
  \item iMessage und RCS zusammen: Ab iOS 18 unterstützt Apple den RCS-Standard. So lassen sich zwischen iPhone und Android hochwertige Bilder, Lesebestätigungen und Tippstatus senden – Schluss mit dem „grünen Blasen“-Bruch. Voraussetzung: Beide Anbieter müssen RCS freigeschaltet haben und die Systemversionen müssen aktuell genug sein.
\end{itemize}

\booksubsection{RCS-Verbreitung in China:}

\begin{itemize}
  \item China Mobile: Frühe RCS-Einführung, in einigen Provinzen bereits „5G Messaging“ mit Rich-Media-Nachrichten und interaktiven Karten. Erfordert ein unterstützendes Endgerät und die Freischaltung im Shop.
  \item China Telecom: Dienst über „Telecom 5G Messaging“. In manchen Tarifen standardmäßig enthalten; Bestandskunden aktivieren ihn per SMS-Befehl.
  \item China Unicom: 5G Messaging in den großen Städten, aber die Geräteabdeckung ist uneinheitlich. Viele Android-Modelle müssen „Erweiterte Nachrichten“ oder „5G Messaging“ manuell in der SMS-App einschalten.
  \item Praxiserfahrung: RCS kommt in China langsamer voran als erwartet – vor allem wegen der starken WeChat-Ökosystem-Dominanz und uneinheitlicher Geräteanpassung. Im Alltag dominiert weiter WeChat; RCS übernimmt eher die Rolle offizieller Push-Nachrichten (Bank, Behörden).
\end{itemize}

\begin{tcbwarn}{Sicherheitsregeln für Bestätigungscodes}
\begin{itemize}
  \item Bestätigungscodes nie weiterleiten: Wer Sie auffordert, „den Code vorzulesen“ oder „die SMS weiterzuleiten“, ist möglicherweise ein Betrüger – selbst wenn er sich als „Kundendienst“, „Paketbote“ oder „Behörde“ ausgibt. Der Code = Ihr Geldbeutel.
  \item Phishing-SMS: Seien Sie misstrauisch bei Links in Nachrichten wie „ETC abgelaufen“, „Versicherungskarte gesperrt“ oder „Paket verloren – Entschädigung“. Seriöse Stellen verlangen nie über SMS-Links Kartennummer oder Passwort.
  \item Nummernerkennung: Offizielle Bank-SMS kommen meist von Kurznummern mit 95 oder von 1069+ mit Unternehmenssignatur. „Bankbenachrichtigungen“ von normalen 11-stelligen Handynummern sind fast immer Betrug.
  \item Zwei-Faktor-Backup: Die Handynummer, die mit wichtigen Konten verknüpft ist, nicht ohne Weiteres wechseln. Wenn es nötig ist, zuerst alle kritischen Konten abmelden – sonst kann der neue Inhaber per SMS-Code Ihre Passwörter zurücksetzen.
\end{itemize}

\end{tcbwarn}

\booksection{WeChat/QQ: Die Super-Apps im Alltag verwalten}

WeChat und QQ sind längst mehr als Messenger: Sie vereinen soziale Funktionen, Bezahlung und Informationszugang in einer „Super-App“.

\booksubsection{Chatverlauf sichern und migrieren:}

\begin{enumerate}
  \item WeChat am PC oder Mac anmelden.
  \item Unten links auf das Menü klicken > [Migration \& Backup] > [Sichern \& Wiederherstellen].
  \item „Chatverlauf auf den PC sichern“ wählen und am Smartphone die gewünschten Chats auswählen.
  \item Migration auf ein neues Handy: Im alten Gerät WeChat [Ich] > [Einstellungen] > [Allgemein] > [Chatverlauf migrieren \& sichern] > [Migrieren] > [Auf Handy/Tablet migrieren], QR-Code scannen – fertig.
\end{enumerate}

\booksubsection{Speicherplatz gezielt freigeben:}

\begin{itemize}
  \item Standardweg: WeChat [Ich] > [Einstellungen] > [Allgemein] > [Speicherplatz]. Hier lassen sich Cache und Chatverlauf nach Zeit oder Dateityp gezielt bereinigen.
  \item Innerhalb eines Chats: Chat öffnen > oben rechts „…“ > [Chatverlauf] > [Speicherplatz verwalten]. So können Sie Bilder, Videos und Dateien einzelner Chats löschen und den Text behalten.
  \item Originale Bilder/Videos: Hochauflösende Originale aus Gruppenchats brauchen am meisten Platz. Unter [Speicherplatz verwalten] „Originale“ markieren und gesammelt löschen.
  \item Automatisches Herunterladen stoppen: [Einstellungen] > [Allgemein] > [Fotos, Videos, Dateien \& Anrufe] und „Automatisch herunterladen“ deaktivieren – so landen nicht ständig große Gruppendateien im Cache.
  \item Mini-Programm-Cache: Unter [Entdecken] > [Mini-Programme] ungenutzte Mini-Programme lange drücken und löschen – das gibt lokalen Speicher frei.
  \item Chatverlauf komprimiert sichern: Beim Sichern in der PC-Version von WeChat „Komprimieren“ aktivieren – mehrere GB Chatverlauf schrumpfen auf 30–50 \% der Originalgröße.
\end{itemize}

\booksubsection{Gruppen-Tipps:}

\begin{itemize}
  \item Gruppenankündigung: Der Admin veröffentlicht sie über „…“ > [Gruppenankündigung]. Alle Mitglieder werden deutlich benachrichtigt – ideal für wichtige Hinweise. Neue Mitglieder sehen sie automatisch.
  \item Gruppen-To-do: Der Admin kann eine angeheftete Nachricht als „Gruppen-To-do“ markieren. Allen Mitgliedern erscheint ein roter Hinweis; nach dem Abhaken verschwindet er. Gut für Anmeldungen und Aufgabensammlung.
  \item Gruppen-Kettenliste: Gib im Eingabefeld „\#Jielong“ ein, um eine Kettenliste zu starten – praktisch zum Sammeln von Bestellungen oder Anmeldungen. Die Ergebnisse werden automatisch oben im Gruppenchat zusammengefasst, damit der Chat nicht zugespamt wird.
  \item Gruppen-Livestream: Oben rechts „+“ > [Gruppen-Livestream]. Startet einen Stream nur für diese Gruppe – ideal für interne Schulungen, Familientreffen oder Produkttests.
  \item Gruppe einklappen: In der Gruppenliste nach links wischen > [Einklappen]. Unwichtige Gruppen landen in „Eingeklappte Gruppen“ und melden sich nur noch bei @Erwähnungen.
  \item Nur Gruppeninhaber/Admins können @alle nutzen: Ab 500 Gruppenmitgliedern können normale Mitglieder nicht mehr alle ansprechen – so wird niemand unnötig gestört.
\end{itemize}

\booksubsection{Moments-Sichtbarkeit fein steuern:}

\begin{itemize}
  \item Gruppen-Sichtbarkeit: Beim Posten „Wer kann das sehen?“ tippen und „Für einige sichtbar“ oder „Für einige verbergen“ wählen. Auf Basis vorheriger Tag-Gruppen (z. B. „Kollegen“, „Familie“, "Kunden") gezielt ausspielen.
  \item Sichtbarkeitszeitraum: [Ich] > [Einstellungen] > [Freundesrechte] > [Moments] > [Zeitraum für Moments]. Wählbar: „Letzte 3 Tage / 1 Monat / 6 Monate / Alle“. Empfehlung: „6 Monate“ als Kompromiss aus Privatsphäre und sozialen Kontakten.
  \item Einzelnen Moment verbergen: Lange auf den Beitrag drücken > [Seine/ihre Moments nicht ansehen] oder [Nicht mehr folgen]. Die andere Person erfährt davon nichts.
  \item Werbung in Moments verwalten: [Ich] > [Einstellungen] > [Über WeChat] > [Persönliche Daten \& Rechte] > [Personalisierte Werbung verwalten]. So lassen sich verhaltensbasierte Empfehlungen abschalten.
  \item Zehn Beiträge für Fremde: [Freundesrechte] > [Moments] > [Fremden zehn Moments anzeigen]. Deaktivieren verhindert, dass Unbekannte Ihre alten Beiträge durchstöbern.
\end{itemize}

\booksection{Videotelefonie im Vergleich}

Jedes Videotelefonie-Tool hat seine Stärken. Das passende Werkzeug für den Anlass macht die Kommunikation reibungsloser.

\begin{itemize}
  \item WeChat-Video: Weit verbreitet, für Jung und Alt. Gut für den Alltag mit Freunden und Familie. Bis zu 20 Personen (über das Mini-Programm „Gruppenvideo“); bei schlechtem Netz sinkt die Bildqualität deutlich.
  \item FaceTime: Exklusiv im Apple-Ökosystem, hervorragende Bildqualität, Ende-zu-Ende-verschlüsselt. Ab iOS 15 können auch Android-/Web-Nutzer per Link beitreten, der Host muss aber ein Apple-Gerät sein. Ideal für Familien mit gemischten Geräten.
  \item Zoom: Erste Wahl für internationale Meetings. Bis zu 100 Teilnehmer, 40 Minuten kostenlos. Bildschirmfreigabe, virtueller Hintergrund und Breakout-Räume sind ausgereift – gut für internationale/Remote-Zusammenarbeit.
  \item Tencent Meeting: Im chinesischen Alltag stark, tief in WeChat integriert, Beitritt per QR-Code. Bis zu 300 Teilnehmer (kostenpflichtig), Cloud-Aufzeichnung, Live-Untertitel und Meeting-Zusammenfassungen – gut für Firmenschulungen und große Meetings.
  \item Empfehlung nach Anlass: 1:1 im Alltag WeChat; unter Apple-Nutzern FaceTime; formelle Business-Meetings Zoom oder Tencent Meeting (in China Tencent, international Zoom).
\end{itemize}

\booksubsection{Bandbreiten- und Datenvolumen-Referenz:}

\begin{itemize}
  \item WeChat-Videotelefonie: ca. 500 Kbps–1 Mbps. Standardqualität ca. 3–5 MB pro Minute, HD ca. 8–12 MB.
  \item FaceTime: SD ca. 1 Mbps, 720p ca. 2 Mbps, 1080p ab 4 Mbps. Ca. 5–15 MB pro Minute.
  \item Zoom/Tencent Meeting: 720p ca. 1,5–3 Mbps, 1080p 3–5 Mbps. Ca. 300–800 MB pro Stunde; der Host verbraucht bei Multi-Personen-Meetings mehr.
  \item Tipp: Lange Videotelefonate möglichst über WLAN. Im Mobilfunknetz die Qualität bewusst senken (in den App-Einstellungen „HD-Video“ deaktivieren), um Datenvolumen zu sparen.
\end{itemize}

\booksubsection{Bildqualität bei Videoanrufen verbessern:}

\begin{itemize}
  \item Licht: Nicht gegen das Fenster, sondern zur Lichtquelle schauen (Tageslicht ist am besten). Gegenlicht lässt das Gesicht dunkel wirken.
  \item Winkel: Die Kamera leicht über Augenhöhe – das vermeidet ein Doppelkinn. Im Querformat das Gerät mit beiden Händen halten.
  \item Audio: Kopfhörer vermeiden Echo und Rückkopplung. In ruhiger Umgebung den Lautsprecher ausschalten.
  \item Netz: WLAN vorziehen. Bei 4G/5G nah ans Fenster gehen, um Empfangsverluste zu vermeiden. In großen Meetings nur die Sprecher-Videos anzeigen lassen.
  \item Hintergrund unscharf: Tencent Meeting und Zoom unterstützen KI-Weichzeichnung – so wirkt auch ein unordentlicher Hintergrund aufgeräumt. Einfarbige Kleidung tragen, damit nicht der Hintergrund „durchscheint“ (Chroma-Key-Effekt).
\end{itemize}

\begin{tcbinfo}{Praxistipp}
Bei Netzproblemen zuerst das Video abschalten und nur weitersprechen – das ist höflicher als sofort aufzulegen. Vor wichtigen Meetings hilft die „Netzwerkprüfung“ von Tencent Meeting, um die Bandbreite vorab zu testen.

\end{tcbinfo}

\booksection{E-Mail: Auf dem Smartphone einrichten und nutzen}

\begin{itemize}
  \item E-Mail-Konto hinzufügen: iOS: Einstellungen > Mail > Konten > Konto hinzufügen. Android hat meist eine „E-Mail“-App.
  \item Praktische Tipps: Unterschiedlichen Konten eigene Benachrichtigungsarten zuweisen. Wischgesten nutzen, um Mails schnell zu archivieren, zu löschen oder zu markieren.
\end{itemize}

\booksubsection{Unterschiede IMAP vs. POP3 vs. Exchange:}

\begin{itemize}
  \item POP3: Mails werden vom Server geladen und dort standardmäßig gelöscht – lokal bleibt nur eine Kopie. Gut für ein einzelnes Gerät und Offline-Lesen. Nachteil: Auf einem neuen Gerät fehlen alte Mails. Wird heute nicht mehr empfohlen.
  \item IMAP: Mails bleiben immer auf dem Server, lokal nur als Cache. Mehrere Geräte (Handy/PC/Tablet) bleiben im gleichen Stand – an einem Ende gelesen, überall gelesen. Das heutige Standardprotokoll.
  \item Exchange (inkl. ActiveSync): Microsoft-Protokoll, synchronisiert neben E-Mails auch Kalender, Kontakte und Aufgaben. Erste Wahl für Firmenmails. Unterstützt u. a. von Outlook, NetEase Enterprise Mail und Tencent Enterprise Mail.
  \item Auswahl: Privatmail IMAP. Firmentmail zuerst Exchange, sonst IMAP. POP3 möglichst meiden, außer es gibt spezielle Archivierungsanforderungen.
\end{itemize}

\booksubsection{Firmenmail einrichten:}

\begin{enumerate}
  \item iOS: [Einstellungen] > [Mail] > [Konten] > [Konto hinzufügen] > „Microsoft Exchange“ oder „Andere“ wählen.
  \item Android: In der E-Mail-App oder Outlook-App „Konto hinzufügen“ wählen, vollständige Adresse und Passwort eingeben.
  \item Serverdaten (Beispiel Exchange): Serveradresse meist „mail.firmendomain.com“ oder „outlook.office365.com“, Domain = Firmenkürzel.
  \item Erscheint „Server kann nicht überprüft werden“, manuell eintragen: Posteingangsserver (IMAP) Port 993, SSL. Postausgangsserver (SMTP) Port 465 oder 587, TLS.
  \item Nach dem Login unter [Kontoeinstellungen] „Push“ aktivieren, um Mails in Echtzeit zu erhalten, und „Mail-Vorschau“ deaktivieren, um Geschäftsgeheimnisse zu schützen.
\end{enumerate}

\begin{tcbinfo}{Kleiner Hinweis}
Firmenmails erfordern oft ein vom Admin freigegebenes „Client-Kennwort“ oder einen „IMAP/SMTP-Authentifizierungscode“ – das Webmail-Passwort allein genügt nicht. Bei Login-Fehlern zuerst im Webmail-Setup einen Code erzeugen und am Handy eingeben.

\end{tcbinfo}

\begin{tcbpractice}{Praxisübung: Einen WeChat-Chatverlauf auf den PC sichern}
\begin{enumerate}
  \item Backup am PC starten: PC-Version von WeChat öffnen > Menü unten links > [Migration \& Backup] > [Sichern \& Wiederherstellen].
  \item Sicherung auf den PC: „Chatverlauf auf den PC sichern“ klicken.
  \item Am Smartphone die Inhalte auswählen und warten, bis die Sicherung abgeschlossen ist.
\end{enumerate}

\end{tcbpractice}
